\subsection{Finite global reduction through the elliptic interfaces}

\begin{statement}{C.2.UN.85}{Proved}{Every reduced SNC torus degeneration has a reversible finite reduction to a classified core}
Let \(f:X\to\Delta\) be a proper degeneration of connected
complex two-tori, with smooth total space and full reduced
SNC central fibre. There is a finite overlap-compatible
modification record over the same disc, isomorphic over
\(\Delta^*\), ending in one of the following actual
smooth-total reduced degeneration germs.

\begin{enumerate}
\item A canonical-flat core. Its central fibre is a smooth
two-torus, one of the exact no-triple elliptic-ruled/diagonal
Hopf cycles of UN65, or one of the marked positive toric
log spaces of UN68 and UN71. The toric alternative has
Euler number equal to its positive triple count.

\item An embedded no-triple core whose components are
minimal fixed-Jacobian elliptic quasi-bundles with identity
integral monodromy, together with any minimal nonlinear
Hopf leaves. All boundary curves are smooth elliptic fibre
supports or the specified Hopf axes. Every elliptic
quasi-bundle base has genus zero or one, and at most one
has genus one. The exact support, including its original
torsors, normal characters, translations and gluing maps,
is one of the rooted-tree or unicyclic supports classified
in UN61 and UN62. Its Euler number is zero.
\end{enumerate}

If every Hopf-birational prime has minimum canonical order,
the first alternative is supplied by UN75. If some such
prime has positive order, the second alternative can be
obtained while retaining a fixed maximal-order minimal
Hopf component and its entire neighbourhood after the
preliminary reduction of UN84. No nonlinear Hopf
neighbourhood is contracted or replaced during this
reduction.

The record retains all actual centres, two-ruling maps,
point and curve contractions, local sigma charts,
normalization branches and restored overlap maps.
Its inverse reconstructs the original analytic family
from this actual endpoint germ. On an originally compact
family all operations take place on that compact base.
If \(S_0\) is the original reduced support, \(S_*\) the
endpoint support, and \(\delta_j\) the signed normal-cone
Euler change of its \(j\)-th retained step as in UN75,
then
\[
                 e(S_0)=e(S_*)-\sum_j\delta_j.
 \tag{C.GP.1}
\]

The classified core central spaces have the compact
occurrence constructions of UN57, UN61, UN62 and UN73.
The inverse record here is applied to its actual endpoint
neighbourhood. Identical central log data in a different
smoothing are not asserted to identify its higher analytic
jets or to carry that particular inverse record.
\end{statement}

\begin{proof}
\par\medskip\noindent\textbf{The initial dichotomy and the protected component.}
Normalize the primitive canonical evaluation as
\[
       \operatorname{div}(\Omega)=\sum_i k_iV_i,
       \qquad \min_i k_i=0.
 \tag{C.GP.2}
\]
If every Hopf-birational prime has order zero, UN75
gives its actual canonical-flat endpoint. UN65 and
UN68--UN71 give exactly the alternatives in (1),
including the actual marked log maps in the positive
case. This uses the proved component restriction that
every class-VII minimal model is Hopf. There is no
unexamined additional class-VII case.

Suppose instead that a Hopf-birational prime has
positive order. Apply the finite truncated record
UN84, including its maximal Hopf minimalization. At
its smooth-total NC endpoint put
\[
                  M=\max_i k_i>0.
 \tag{C.GP.3}
\]
The minimum is still zero and all orders are at
most this fixed \(M\). Choose one of the surviving
minimal Hopf components \(H\) of order \(M\).
UN82 shows that it has exactly one elliptic
boundary curve \(J\), with torsion normal, and
that its neighbour has smaller order. The marked
neighbourhood is embedded. No elliptic self-double
curve has been introduced in the preliminary
record: its new incidences are on rational
exceptional curves, as proved in UN84.

All later whole-component removal blocks will be
the positive-order blocks of UN77 and UN81. Such
a ruled component of order \(q\le M\) has one
horizontal boundary section, whose neighbour has
order \(q-1\), and every other boundary is rational
and vertical. It cannot meet \(H\): a rational
curve cannot lie on a minimal Hopf, and the
section case would require \(q=M+1\). A plane
contraction also has only rational boundary.
The same argument applies to the cleanup blocks
and their missing-fibre insertions. Direct moves
below will likewise avoid \(H\). Thus its entire
neighbourhood is protected and provides the same
starting interface after every completed block.

\par\medskip\noindent\textbf{The recorded models, not a numerical eligibility class.}
At a smooth-total NC stage, call a component
\emph{recognized} only when one of the actual
comparisons of UN83 has been recorded for its
smooth normalization \(T\), together with its
entrance elliptic curve and its point map
\[
                         \pi_T:T\longrightarrow Z_T.
 \tag{C.GP.4}
\]
There are three recorded types: an elliptic
quasi-bundle model; a nonlinear Hopf model with
its one axis; or a one-section elliptic ruled
model with its actual rank-two bundle extension.
The point map is an isomorphism near the entrance
curve. All boundary branches, including both
branches of a self-double, are retained.

Let \(d_T\) be the exact number of point blowups
in this map. The following properties are part
of those actual comparisons:
\[
\begin{gathered}
 d_T\ge0,\qquad
 \operatorname{Exc}(\pi_T)\subset D_T,\\
 \text{every rational boundary is exceptional
 and has negative square},\\
 \text{every elliptic boundary has nonpositive square}.
\end{gathered}
 \tag{C.GP.5}
\]
For the elliptic type, the boundary is vertical
for the composed elliptic map; (C.GP.5) follows
from UN80 and the fibre squares. For the other
two types the entrance curve has square zero,
all point centres are away from it, and every
other boundary is exactly a positive discrepancy
exceptional curve, by (C.TI.2)--(C.TI.3).

Initially recognize \(H\) alone. UN83 applied
to its actual isolated torsion-normal axis gives
the elliptic or nonlinear Hopf type, and its
point defect is zero. Let
\[
\begin{gathered}
 u=\#\{\text{current unrecognized global component primes}\},
 \\
 d=\sum_{T\ \mathrm{recognized}}d_T.
\end{gathered}
 \tag{C.GP.6}
\]
Recognition changes no complex space: it records
the actual maps of UN83 and their marked curves.
In particular the word does not assert that an
unrecognized receiving surface can never acquire
such a comparison at a later step.

\par\medskip\noindent\textbf{A direct move on a recorded elliptic model.}
Suppose that a recognized elliptic-type component
has positive defect. Choose a last exceptional
curve of (C.GP.4). It is an actual boundary by
UN80. The complete elliptic fibre graph and the
exact receiving-degree comparison (C.AP.5)--
(C.AP.6) apply without any additional sign
assumption on the entrance coefficient.

Here is the exact comparison with an embedded
curve in the total space at an NC stage. On a
complete fibre graph of the point map, write
\(a_C\) for the coefficient of each component
in the actual canonical divisor on the
normalization, assigning coefficient zero
also to a component absent from the actual
boundary. Adjacent shifted coefficients
\(a_C+1\) are relatively prime. There is no
initial edge on the smooth elliptic fibre;
a blowup at a smooth point creates consecutive
shifted coefficients, and a node blowup inserts
the sum of the two old shifted coefficients.
Thus the assertion follows through the actual
point map, including negative initial coefficients.
In particular no edge has coefficients
\((-1,-1)\).

A self-double branch has coefficient \(-1\).
If the image of the chosen last \((-1)\)-curve
were nodal at an actual triple point with two
local branches of its source component, its
two neighbours on the normalization would be
the two self-double branches. The last-node
coefficient formula would also give
\(a_E=(-1)+(-1)+1=-1\), contradicting that
edge invariant. A triple with three branches
of the source gives adjacent \(-1\) branches
directly. With only one local source branch,
two image branches would already make a node
on the chosen curve in that smooth source
normalization, contrary to its being the
smooth last exceptional \(\PP^1\).
The actual image curve is consequently
smooth. For a self-double its normalized
degree-two cover is unramified, by the local
NC coordinates. A connected such cover by
\(\PP^1\) is impossible by Riemann--Hurwitz;
the cover therefore splits into two distinct
normalized branches. Over the resulting
simply connected image \(\PP^1\), those
labels give actual two smooth surface germs.
When the triple count is two, the normal-product
map gives the two normal degrees \(-1,-1\).
The actual type-II neighbourhood and its two
ruling maps therefore apply even if these
two surface germs belong to one global prime.
For the positive-discrepancy models below,
all rational coefficients are positive, so
the same self-double obstruction is immediate.

With two triple points, the actual type-II move
blows down this curve on its two adjacent
normalized surfaces. Each opposite receiving
boundary has square at least zero, or at least
one when it is elliptic. By (C.GP.5) neither
receiver was recognized before this move. Thus
the move adds no point defect to a recognized
surface.

If its second source surface was recognized,
the contracted rational boundary was exceptional
for its recorded point map, again by (C.GP.5).
That map factors through the actual blowdown:
it is constant on the contracted curve, and
the proper surface contraction and its direct
image of the structure sheaf give this descent.
The resulting birational surface morphism is
the remaining sequence of point blowups. Its
defect decreases by one. The same factorization
applies to the first source. All recorded
entrance curves are disjoint from these
exceptional curves and are unchanged. Hence
the recognized set and its models persist,
\(u\) is unchanged, and \(d\) strictly decreases.
If both source germs belong to one global
recognized component, their two normalized
curves are disjoint. The same point map
factors through both ordinary contractions;
its defect drops by two while the global
component is counted only once.

If a last curve has one triple point, use the
actual positive ruled-removal block of (C.AP.7).
If it has none, use the exact plane progress
block (C.AP.8)--(C.AP.13). Each gives a strict
component-card decrease after its finite cleanup.
No nonboundary last-point case is admitted: its
exclusion is the actual field statement UN80,
not an inference from nonzero canonical coefficients.

\par\medskip\noindent\textbf{A direct move on either positive-discrepancy model.}
Now let the recognized model be nonlinear Hopf
or one-section ruled, with positive defect.
The actual canonical divisor is
\[
 K_T=-aJ_T+\sum_E b_EE,
 \qquad b_E>0,
 \qquad a\ge2,
 \tag{C.GP.7}
\]
with \(a=2\) for the ruled type. Every point
centre is away from the preserved \(J_T\),
and the full other boundary is exactly the
exceptional set. The graph of each exceptional
tree starts with a rational curve of square
\(-1\) and valence zero. If \(C^2=-b\) and
\(n\) is its valence in this entire exceptional
set, the actual smooth and node blowups give
\[
                              n\le b+1.
 \tag{C.GP.8}
\]
At a smooth point both \(n\) and \(b\) increase
by one on an old component; at a node \(b\)
increases and its old valence is unchanged.
A newly inserted curve has \((b,n)=(1,0)\),
\((1,1)\), or \((1,2)\), according as its
centre is off the exceptional support, on
one component, or at a node. These are
exactly the updates proving (C.GP.8).
At the last node the opposite
square is \(b+1-t\ge0\), since the actual
triple count \(t\) is at most the old valence.
All opposite curves here are rational.
Thus no receiver is already recognized by
(C.GP.5), and the exact point-map factorization
in the preceding paragraph again strictly
decreases \(d\). The positive discrepancies
also give the precise height comparison:
if the two old coefficients are \(b_1,b_2>0\)
and the current component has order \(p\),
the receiving orders are
\[
                 p+b_1+1,\qquad p+b_2+1.
 \tag{C.GP.9}
\]

A last curve with no triple point has a
rational \(+1\) receiver of order \(p+2\),
and (C.AP.8)--(C.AP.13) supply the finite
plane progress block. With one triple point
its receiver is ruled with one section and
has positive order, so the same ruled block
applies. These changes concern only the
actual rational exceptional configurations,
not the nonlinear Hopf neighbourhood itself.

When a one-section ruled model reaches defect
zero, it is the actual bundle \(\PP(V)\)
with sole boundary section \(J_T\), including
its possibly nonzero extension class. UN74
gives the ordinary curve contraction along
its ruling. Its order is
\(k_T=k_{\mathrm{neighbour}}+1>0\), by UN83.
This is another whole-card decrease.

\par\medskip\noindent\textbf{Recognition across the next actual interface.}
We always eliminate positive recorded defects
before recognizing another component. We also
remove a defect-zero one-section ruled model
by the preceding actual contraction. If neither
action is available, every recognized component
is a minimal elliptic quasi-bundle or a minimal
nonlinear Hopf leaf.

Such a component is embedded and has disjoint
smooth elliptic boundary curves. Indeed no
rational boundary can remain when its point
defect is zero. The original fibre was SNC,
and the retained type-II and one-section
sigma records create new self-incidences only
on rational fibre or exceptional curves.
Their generic one-section contraction cannot
identify two elliptic horizontal branches.
Thus there is no new elliptic self-double
curve on one of these minimal components.

If an unrecognized component remains, the
connected incidence graph of the central
support has an edge from the recognized
region to its complement. Starting with
\(H\) and recording each new entrance edge
makes that region connected; the direct
moves have retained these entrance curves.
The edge is therefore a smooth elliptic
curve, isolated on its recognized side.
There is no triple point on it, so its
other normalization branch is isolated
as well. Its normal is torsion: on a
quasi-bundle it is the exact normal root
of a smooth multiple fibre, and on a
nonlinear Hopf it is its axis character.
The actual reduced-fibre normal-product
map gives the inverse torsion line on
the other side.

Apply UN83 to that actual other normalization.
It supplies one of the three point models
(C.GP.4), with the entrance neighbourhood
unchanged and every original boundary
retained. Record it as recognized. Then
\(u\) decreases by one. The newly recorded
defect may be any finite nonnegative integer;
we make no assertion that this recognition
alone decreases \(d\).

\par\medskip\noindent\textbf{Whole blocks may reset the recognized region.}
A whole-component removal and its sigma
cleanup need not preserve the point models
on the previously recognized region. For
example a missing-fibre insertion can create
additional point blowups on a section parent.
We allow all such changes explicitly.
After this completed block discard every
recognition except the unchanged maximal
Hopf \(H\), and restart with its same
torsion-normal interface. This is possible
because its entire neighbourhood is protected
as proved above. There is no assertion that
the rest of the processed region was protected.

Let \(c_j\), for \(1\le j\le M\), count
the retained global components of order \(j\).
The full fixed potential is
\[
                 (c_1,c_2,\ldots,c_M,u,d).
 \tag{C.GP.10}
\]
Every direct type-II move preserves the card
vector and \(u\), and strictly decreases \(d\).
Recognition preserves the cards and strictly
decreases \(u\), regardless of its finite
new defect. Every whole-component block
strictly decreases its first changed card,
regardless of the finite reset of \(u,d\).
The actual missing-fibre and cleanup rules
add only finitely many cards at the next
order and add none at \(M\). Thus all
retained orders stay at most the fixed \(M\).

Lexicographic order on this fixed finite
product of nonnegative integers is well
founded. The internal plane progress and
sigma cleanup blocks are finite by UN81
and UN77. Consequently this procedure
cannot continue indefinitely.

At its endpoint no recorded defect is
positive, no one-section ruled model
remains, and no component is unrecognized.
Otherwise the preceding paragraphs give
another strict decrease. All components
are therefore the minimal elliptic
quasi-bundles and nonlinear Hopf leaves
in (2), and every boundary is smooth
elliptic and isolated. In particular the
endpoint is embedded SNC and has no
triple point.

\par\medskip\noindent\textbf{Identification of the endpoint and the retained analytic scope.}
Every endpoint component and double curve
has topological Euler number zero, so
\(e(S_*)=0\). The exact invariant relative
cohomology frame gives
\(p_g(Z)=g(B)\) for a quasi-bundle. The
common regular Jacobian is retained by
the actual normal-root maps: along an
edge joining two quasi-bundle vertices,
the dual normal lines have the same
exact order, and inversion identifies
their connected root covers of the
elliptic double curve. These are exactly
the two regular elliptic fibres, by
(C.QT.2)--(C.QT.3); their translations
are retained in the gluing record.
At a nonlinear Hopf leaf the same root
cover is instead its exact period cover
of UN50. Such leaves do not interrupt
propagation through the connected
quasi-bundle region, and their actual
leaf periods and finite characters are
retained by that comparison.
The
compact-core bound UN28 therefore permits
at most one genus-one base. If there is
one, that normalization has \(p_g=1\).
UN47 gives the actual surjection from its
fundamental group onto that of the full
support. The graph-of-spaces comparison
gives a surjection from the support
fundamental group to that of its incidence
graph, and kills the image of this
component's fundamental group.
The graph fundamental group is therefore
trivial, so the connected graph is a tree.
The exact coherent restriction and
fundamental-group comparisons of UN61
then give its rooted-tree support, with the
actual integral Chern and recursive
period conditions. If there is none,
UN62 gives the unicyclic support and
its exact torsor, height and common-period
conditions. Its nonlinear leaves retain
the secondary finite characters of UN50.
These applications keep the original
surfaces and gluing maps; they do not
replace the actual endpoint smoothing.

All prior operations are the actual
finite diagrams already used in UN75,
UN77, UN81 and UN84. Their proper graphs
glue by the retained overlap maps, so
their inverse recovers this given
analytic family. Equation (C.GP.1)
is the telescoping normal-cone formula
on that finite record, including its
restored-overlap subdivision. Annular
gluing to an unchanged compact outside
gives the same actual modifications
when the starting total space is compact.

The separate compact constructions for
the classified core support data prove
their occurrence. They do not identify
arbitrary higher transverse jets of
this actual endpoint with another core
smoothing. No inverse sigma record is
transferred through an unproved such
identification. This is the claimed
finite reduction and exact reconstruction
of every given reduced SNC torus
degeneration.
\end{proof}
